% ============================================================
% !TeX root = ../main.tex
%  Abstract (English)
% ============================================================
\chapter*{Abstract}

\noindent
Short-axis cardiac magnetic resonance imaging shows the left ventricle
through separate slices, so the regions between them must be reconstructed to
obtain a continuous representation of the ventricular wall. Interpolation
algorithms can connect the available contours, but they do not guarantee that
the resulting shape corresponds to the patient's anatomy. Regional clinical
assessment uses the American Heart Association (AHA) 17-segment model, yet
visual assessment remains observer-dependent, and thickness measurements on
selected slices do not describe the entire ventricular wall.

This thesis develops a machine-learning approach that reconstructs the left
ventricle from these contours. An implicit neural representation describes
the inner and outer surfaces of the ventricular wall continuously and
considers cardiac phase, allowing the ventricle to be reconstructed at
end-diastole and end-systole. Training begins with synthetic ventricular
shapes, followed by fine-tuning on data derived from clinical segmentations.
Wall thickness is computed at local points on the reconstructed surface and
summarised within the seventeen AHA segments.

The reconstructions are compared with methods based on radial basis functions
and a statistical shape model. The reconstructed surfaces follow the observed
contours, but the methods agree more closely on the cavity than on the
myocardial wall. Thickness measurements depend on both the reconstructed shape
and the algorithm used to calculate them, so a strong correlation between
measurements does not mean that their values are equivalent. The AHA-based
analysis also captures wall thickening between cardiac phases.

The work links continuous ventricular reconstruction to objective,
geometry-based wall-thickness measurements at local and AHA-segment levels.
Without an independent three-dimensional anatomical reference, it is not
possible to confirm which reconstruction best represents the true anatomy,
but the results show how the choices made to complete the ventricular shape
influence the analysis of its thickness.

\vspace{1em}
\noindent\textbf{Keywords:} Cardiac magnetic resonance imaging; left ventricle;
three-dimensional reconstruction; implicit neural representations; signed
distance functions; myocardial wall thickness.

\clearpage
